\iflanguage{english}
{\chapter{Fazit und Ausblick}}
{\chapter{Conclusion and Further Work}}

\label{sec:conclusion}

Through an extensive series of experiments, I systematically evaluated the impact of various super-resolution models on the quality of UV texture maps. Among them, RCAN performed best in the $256 \rightarrow 512$ super-resolution task, SRFormer demonstrated superior performance in the $512 \rightarrow 1024$ super-resolution task in the UV domain, achieving high pixel-level fidelity.

 In me experiments, although many HRHumanTex variants perform better than the SMPLitex baseline, HRHumanTex-SwinIR and HRHumanTex-DiffBIR consistently exhibited strong and stable performance across both UV-space evaluation and inpainting-based reconstruction evaluation in all variants. These models not only produced perceptually plausible textures in UV space but also preserved details and consistency in rendered 3D views.

The final HRHumanTex framework, built upon SwinIR and DiffBIR, achieves stable high-resolution human texture estimation from a single RGB image, demonstrating robustness across diverse identities and poses.

Nevertheless, my work still faces several limitations. First, when the subject appears in profile views, the resulting facial texture often lacks symmetry or completeness. Second, large garments like oversized shirts may cause texture bleeding, where shirt textures mistakenly extend into pants regions due to partial observation ambiguity. I will show some failure cases and special cases in the Appendix.

In future work, I plan to improve the differentiable rendering-based texture refinement module to further enhance spatial coherence and semantic correctness. Such optimization can explicitly enforce consistency between the rendered appearance and the input image, especially in occluded or ambiguous areas, potentially improving both realism and generalization.

